\documentclass[11pt,a4paper]{article}

\usepackage[utf8]{inputenc}
\usepackage[cyr]{aeguill}
\usepackage[francais]{babel}

% Les packages
%\usepackage[french]{babel}
\usepackage{fancybox}
\usepackage{graphics}
\usepackage{color}
\usepackage{latexsym}
\usepackage{amsmath,amsthm,amssymb}
\usepackage[plain]{fullpage} 
\usepackage{verbatim}
\usepackage{times,epsfig}
\usepackage{listings}
\usepackage{multirow}
\usepackage{amsmath,amsthm,amssymb,amscd,txfonts,bbm} % RAJOUT ICI !!!
\usepackage{graphics,epsfig} % RAJOUT ICI !!!
%\usepackage{algorithmic}
\usepackage[final]{pdfpages} 
\usepackage{fancyvrb}
\usepackage{enumerate}
\usepackage{listingsutf8}
\usepackage{comment}
\usepackage{minted}

% small et verbatim: to be coherent with \file
\newenvironment{smallverbatim}%
{\endgraf\small\verbatim}{\endverbatim}

\newenvironment{qv}
{\quote\Verbatim}
{\endquote\endVerbatim}

\usepackage{lastpage}

\newcommand{\terme}[1]{\texttt{#1}}
\newcommand{\termSet}[1]{\{\texttt{#1}\}}
\newcommand{\guill}[1]{<<\,#1\,>>}
% Macro pour l'inclusion de figure  
\newcommand{\fig}[2]{%
  \begin{figure*}[hbt]
   \begin{center}
      \includegraphics[width=6cm]{#1}
  \caption{\label{#1}  #2}
  \end{center}
 \end{figure*}
}
\newcommand{\figTex}[2]{%
  \begin{figure*}[hbt]
 \centerline{\input{#1.pstex_t}}
  \caption{\label{#1}  #2}
 \end{figure*}
}

% Macro commentaires
\definecolor{turquoise}{rgb}{0.20 ,0.60, 0.60}
\definecolor{vert}{rgb}{0.40,  0.60 ,0.00}
\definecolor{violet}{rgb}{0.80 , 0.00 , 0.80}
\definecolor{bleu}{rgb}{0.20,  0.20,  0.80}
\definecolor{rose}{rgb}{1.00,  0.20,  0.40}

\lstset{basicstyle=\footnotesize,showstringspaces=false,language=XML,breaklines=false,columns=fullflexible,frame=shadowbox, captionpos=b}

\newenvironment{solution}[0]{
~\\ \color{red} \textbf{Solution : }\color{black}}{%
}
\excludecomment{solution}


\parindent=0pt           
                             
\usepackage{url}

\newtheorem{Exercice}{Definition}


\begin{document}

\title{Sujet d'examen UE NFE204}
\author{Session 1 PROBTP}
\date{30 mars  2018}

\maketitle

\begin{center}
{\bf \Large Consignes}\\
\begin{tabular}{| p{15cm} |}
\hline
Dur\'ee de l'examen : 3h00; Documents non autoris\'es;  Calculatrice autoris\'ee\\

Les exercices sont indépendants les uns des autres. Merci de soigner votre écriture.\\
Ce sujet contient \pageref{LastPage} pages.\\

\textit{Important}\,: nous n'évaluons pas les réponses par la syntaxe. Ne vous
inquiétez pas d'utiliser un point-virgule à la place d'un point d'exclamation ou
autres détails mineurs
\\
\hline
\end{tabular}
\end{center}

\emph{L'ensemble de l'examen doit vous permettre de montrer que vous avez assimilé les principales notions
du cours et que vous savez les mettre en application et les expliquer. La clarté, la concision,
la précision des réponses sera donc appréciée.}


\section*{Première partie : modélisation semi-structurée (6 pts)}

Un grand magasin enregistre tous les tickets de caisse des transactions effectuées par ses clients.
Un ticket de caisse comprend la date et l'heure de chaque transaction, le numéro de la caisse, le prix total,
et la liste des articles achetés avec, pour chacun, la quantité et le prix unitaire.  Le ticket de caisse comprend également
l'identifiant du client si ce dernier a utilisé sa carte de fidélité. Sinon cet identifiant
est à \texttt{null}.


\begin{enumerate}
  \item Dans une base relationnelle, combien de tables faudrait-il pour représenter 
         les informations d'un ticket de caisse (expliquer et donner le schéma de ces tables); combien d'insertions de ligne
             faut-il faire pour chaque ticket\,?
    \item Si on choisit une base NoSQL, comment représenter un ticket de caisse en JSON.
        Donnez un exemple simple. Combien d'insertions faut-il faire pour crér un ticket\,?
   \item Je veux compter le nombre de transactions par caisse et par jour.  Quelle est la requête SQL si les données
       sont dans une base relationnelle\,?
    \item Même question si la base est NoSQL\,: comment calculer le nombre de transactions par caisse et par jour
    \item Indiquez les avantages et inconvénients de chaque solution, relationnelle ou NoSQL.
 \end{enumerate}


\begin{solution}

\begin{enumerate}
   \item En relationnel, il faut 2 tables, une pour les tickets avec les champs \texttt{noCaisse},
      \texttt{date}, \texttt{heure} et \texttt{prixTotal}, une autre pour les articles avec \texttt{idArticle},
         \texttt{quantité} et \texttt{prixUnitaire}. Un champ \texttt{idClient}
          dans la table \texttt{Ticket} peut être à \texttt{null}. La table \texttt{Article} comprend
         une clé étrangère \texttt{idTicket} pour savoir de quel ticket l'achat de l'article fait partie.
         
          Il faut une insertion dans la table \texttt{Ticket},
           et autant d'insertions dans la table \texttt{Article} qu'il y a d'articles achetés dans
           une même transaction. 
      
      \item En JSON, par exemple:

\begin{minted}[frame=leftline,
               baselinestretch=.9,
               fontsize=\footnotesize,
               framesep=2mm]{json}
{
  "_id": 978,
  "noCaisse": 12,
  "date": "01/01/2017",
  "heure": "13:30",
  "prixTotal": 58.47,
  "idClient": null,
  "articles": [
     {"idArticle": 1982, "quantité": 5,   "prix": 10.45},
     {"idArticle": 34, "quantité": 1,   "prix": 45},
     {"idArticle": 154, "quantité": 3,   "prix": 3.02},
  ]
}
\end{minted}

Une insertion suffit.

\item En SQL:

\begin{minted}{sql}
select count(*) from Ticket t, Article a 
where a.idTicket = t.id
group by date, noCaisse
\end{minted}

\item En NoSQL, il faut écrire et exécuter un programme MapReduce.

La fonction de Map émet  une paire intermédiaire dont la clé est le critère de regroupement,
soit la paire (date, numéro de caisse).

\begin{minted}{bash}
function fonctionMap ($ticket)
{
  emit ( [$ticket.date, $ticket.noCaisse], 1)
}
\end{minted}

La fonction de Reduce effectue le compte des paires intermédiaires, autrement dit des tickets
ayant même  date et même numéro de caisse.

\begin{minted}{bash}
function fonctionReduce ($clé, $valeurs)
{
  return ($clé, count($valeurs))
}
\end{minted}
\end{enumerate}
\end{solution}


\section*{Deuxième partie : MapReduce (3 pts)}

Je veux calculer le montant moyen des tickets par client, pour les clients dont
 l'identifiant est non nul. 

Décrivez la modélisation MapReduce de ce calcul. 
 Donnez le pseudo-code de chaque
fonction, ou indiquez par un texte clair son déroulement.


\section*{Troisième partie : comprendre ElasticSearch (8 pts)}

Voici quelques aspects d'Elastic Search qui devraient vous être familiers, 
et d'autres avancés qui n'ont pas été abordés en cours

\subsection*{Indexation}

La documentation nous dit\,:
\textit{When you index a document, it is stored on a single primary shard determined by a simple formula\,:}

    $$shard = hash(id) \%\ \rm{number\_of\_primary\_shards}$$
    
Voici la configuration de votre \textit{cluster} ElasticSearch pour les questions qui suivent.
Nb\,: réplica = $n$ signifie $n+1$ copie d'un document.

   \begin{smallverbatim}
cluster.name: philippe
index.number_of_shards: 3
index.number_of_replicas: 2
   \end{smallverbatim} 

Et je crée un index. Chacune des questions ci-dessous vaut 1 point. Une brève explication sera appréciée.

\begin{enumerate}
       \item Est-ce que mon \textit{cluster} peut être constitué d'un seul n{\oe}ud\,?
       \item Est-ce que je peux étendre le nombre de \textit{shards}\,?
       \item Est-ce que je peux étendre le nombre de réplicas\,?
       \item Est-ce que je peux ajouter des n{\oe}uds à mon \textit{cluster}\,?
       \item Si oui, à partir de quand est-ce que ça ne sert plus à rien d'ajouter des n{\oe}uds\,?
    \end{enumerate}

Un dernier point à gagner. Un de vos collègues vous affirme qu'il vaut mieux
créer un index avec beaucoup de \textit{shards}, pour anticiper la croissance future. Voici
un argument contre ce choix, issu de la documentation\,:
\textit{term statistics, used to calculate relevance, are per shard. Having a small amount of 
data in many shards leads to poor relevance.} 

Comment lui expliqueriez-vous cet argument dans un français clair et concis\,?

\subsection*{Requêtes}

\begin{enumerate}

  \item On effectue une recherche dans un \textit{cluster} ElasticSearch constitué de 10 fragments (primaires), et on veut récupérer
      le résultat par tranches de 10 documents. Combien faut-il récupérer au minimum de documents de chaaue 
      fragment pour constituer un résultat correct pour chaque tranche\,? 

  \item Quelle information le coordinateur de la requête
    devrait-il propager à tous les fragments pour que le classement soit correct (oui cette question est liée
    à un extrait de la doc donné précédemment)\,? 

 \end{enumerate}



\section*{Quatrième partie : questions de cours (3 pts)}


\begin{enumerate}
  \item En recherche d’information, que signifient les termes ``faux positifs'' et ``faux négatifs''?
  \item Rappeler la règle du quorum (majorité des votants) en cas de partitionnement de réseau, et justifiez-la.
   \item Qu’entend-on par ``bag of words'' pour le modèle des documents textuels en recherche d’information?
\end{enumerate}


\end{document}
